\begin{titlepage}
\centering
{\Large Notas de entrevista técnica\par}
\vspace{1.0cm}
{\huge\bfseries Entrevista con Hong Letong (洪乐潼): AI for Math, Lean e intuición matemática\par}
\vspace{0.6cm}
{\Large Zhang Xiaojun conversa con Hong Letong (Carina Hong)\par}
\vspace{0.4cm}
{\large Elaborado a partir del video público de Zhang Xiaojun Podcast, una transcripción generada con GPU y diagramas conceptuales propios\par}
\vspace{0.3cm}
{\large 2026-04-20\par}
\vspace{0.9cm}
\includegraphics[width=0.88\textwidth,height=0.42\textheight,keepaspectratio]{cover.jpg}\par
\vfill
\begin{tcolorbox}[width=0.92\textwidth, colback=black!2!white, colframe=black!55, sharp corners]
\textbf{Título del video}: 137. Entrevista de 4 horas con Hong Letong: AI for Math, convertir las matemáticas a Lean, las demostraciones de «El Libro» de las matemáticas, la intuición, lo creado y lo descubierto\par
\textbf{Canal del video}: Zhang Xiaojun Podcast\par
\textbf{Duración del video}: 04:24:17\par
\textbf{Enlace del video}: \href{\videourl}{\nolinkurl{\videourl}}\par
\textbf{Descripción del material}: YouTube no ofrece subtítulos; el cuerpo del texto se elaboró a partir de la transcripción con faster-whisper large-v3 en CUDA, los capítulos del video, la portada y figuras didácticas propias. Las tomas fijas de la entrevista no se reutilizan en el cuerpo del texto.\par
\textbf{Página del proyecto}: \href{\repourl}{\nolinkurl{\repourl}}\par
\end{tcolorbox}
\end{titlepage}

\tableofcontents
\newpage
